\subsection{(a) Generic separable disutility}

Instead of a constant marginal utility per kWh consumed, the consumer now has a preferred
consumption profile $\lref$ (kWh/h) and experiences a disutility $D_t(\load)$ (DKK) when its
load deviates from it. $D_t$ is convex, differentiable, non-negative, and equal to $0$ at
$\load = \lref$. All the other parameters, variables and constraints are those of
Question~1, and the disutility is still separable per hour.

\subsubsection{Objective function}
The term $\uL \, \load$ of \eqref{eq:daily_net_utility} is replaced by $-D_t(\load)$:
\begin{equation}
    \max_{\{\load, \qimp, \qexp, \pvprod\}_{t \in \T}} \quad
    \sum_{t \in \T} \Big( - D_t(\load) - \pimp \, \qimp + \pexp \, \qexp - \cpv \, \pvprod \Big)
    \label{eq:q2_objective}
\end{equation}
subject to the same constraints \eqref{eq:c_balance}--\eqref{eq:c_export_nonneg} as in Question~1.

\subsubsection{Properties compared with Question 1}
\begin{itemize}
    \item \textbf{Problem class and convexity:} since $D_t$ is convex, $-D_t$ is concave, so
    \eqref{eq:q2_objective} maximizes a concave function over the same polyhedron as in
    Question~1: it is still a convex optimization problem, but in general no longer an LP. It
    stays an LP only if $D_t$ is piecewise linear, and becomes a QP if $D_t$ is quadratic.
    \item \textbf{Decomposability:} it still holds. $D_t$ only involves the load $\load$ of
    hour $t$, so every term of the objective and every constraint still involves a single
    hour, and the hourly decomposition argument of Question~1 applies unchanged.
    \item \todo{Feasibility compared with Question 1.}
\end{itemize}

\subsubsection{Question 1 as a particular case}
Choosing the constant reference $\lref = \Lmax$ and the linear disutility
\begin{equation}
    D_t(\load) = \uL \, (\Lmax - \load)
\end{equation}
gives $-D_t(\load) = \uL \, \load - \uL \, \Lmax$, that is, the objective of Question~1 up to
the constant $-24 \, \uL \, \Lmax$, which does not change the optimal solution. On the feasible
set $\load \leq \Lmax$, so $D_t \geq 0$, and $D_t = 0$ exactly at the reference.
\todo{Which of the assumptions on the disutility (convex, differentiable, non-negative, zero
at the reference) does this choice relax? Check it against the brief.}
